\section{Toward Integrated, Targeted, and Adoptable Solutions}

The preceding frontiers make clear that better process understanding, models,
and monitoring are necessary but not sufficient: translating frontier science
into water-quality and productivity gains depends on how practices are
combined, where they are placed, how they perform under a changing climate, and
whether they are economically and institutionally adoptable. This section draws
the integrative and socio-economic strands of the literature together. Its
recurring message is that no single practice suffices, that biophysical
questions are inseparable from human ones, and that the decisive advances will
come from coupling mechanistic understanding with spatial targeting, economics,
policy, and shared decision-support infrastructure.

\subsection{Bundling, sequencing, and spatial targeting of practices}

A pervasive conclusion is that no single best management practice (BMP) is a
silver bullet \parencite{Sunohara2016}, and that multiple BMPs in series are
needed---with in-field/edge-of-field combinations judged highly promising yet
little studied, and wetland and water-retention measures still to be developed
for sensitive downstream waters \parencite{Lenhart2016}. Attaining water-quality
goals will require combinations implemented at a high level in a targeted manner
\parencite{Christianson2014}. Open questions concern how best to pair
conservation tillage with cover crops and in-field with edge-of-field practices
while accounting for legacy effects \parencite{Makhtoumi2025}, how to formalise
optimisation of synergistic combinations \parencite{Kreig2019}, and how
practices interact and what nitrogen-removal capacity each provides given better
mapping of tile-drained land \parencite{McLellan2015}. Broad programmes are
proposed to compare drained and undrained systems, to bundle bioreactors,
phosphorus reactors, controlled drainage (CD), and low-impact development
alongside internal-loading management \parencite{Baulch2019,Husk2024,Husk2017},
to develop cost-effective BMPs that support cover crops and mitigate warming
\parencite{Gupta2023}, and to test intercrops and hydrology-altering management
for phosphorus reduction \parencite{HernandezRamirez2011,GAYNOR1995b}.

Closely related is targeting practices to the parts of the landscape that matter
most. Identifying hydrologically active runoff-generating areas as mitigation
priorities is recommended \parencite{Melland2008}, but whether elevated loads
originate from many diffuse sources or a few hot spots
\parencite{Irvine2019}, whether targeted or regionally saturated cover-crop
implementation is needed for consistent reductions \parencite{Speir2022}, and how
to delineate zones of higher nitrate leaching \parencite{Bakhsh2005b} remain
open. Targeting critical source areas and hydrologically connected hotspots is
repeatedly stressed \parencite{McDowell2015,Brendel2019}, prioritisation is
expected to improve the return on conservation investment
\parencite{Makhtoumi2025}, and automating spatial-targeting principles into an
accessible GIS tool with whole-basin surveys and cost--benefit analysis is a
priority \parencite{Kvítek2023}. Further scaling needs include combining GIS with
hydrological modelling and saturation risk to find nitrate thresholds across
soil texture and temperature \parencite{Smith2023}, resolving how drainage
density controls stream solute dynamics \parencite{Kennedy2012}, and identifying
critical precipitation periods rather than annual averages \parencite{Salo2006}.

\subsection{Climate change, extremes, and system resilience}

Climate adaptation frames many priorities. The overarching need is to keep
identifying practices and combinations that increase system resilience to
climate extremes across additional soils and longer periods
\parencite{Saha2024}, supported by modelling agendas spanning additional
periods, representative concentration pathways, management options, cultivars,
cover-crop types, and soils \parencite{Malone2020}, and by field-scale
integrated biophysical modelling that resolves within-county variation in
fertiliser and manure inputs under multiple CMIP5/CMIP6 forcings with attention
to perennial-grass co-benefits \parencite{Ferin2023}. A central design
challenge is cropping and drainage systems that function under both drought and
flood, including engineering solutions to capture excess floodwater
\parencite{Dietzel2016}. Adapting drainage design itself to a changing climate is
an explicit priority, tied to understanding weather--soil--groundwater--discharge
interactions under increased precipitation and to evaluating designs over longer
periods and different climates
\parencite{Deelstra2015,Tuohy2018,Malota2024,Deelstra2014,Deelstra2014b}.
Remaining needs span the costs, benefits, and payback of adaptation and
carbon-farming knowledge exchange \parencite{PeltonenSainio2021}, the effect of
extreme events on cover-crop performance \parencite{Gupta2023b}, nitrogen
consequences under warming \parencite{Haghnazari2020}, and testing across
additional rainfall environments, soils, and rotations to confirm
drainage--management interactions
\parencite{Takeshima2023,Liu2025,Grewal2025,Negm2014,BENOIT1967}.

\subsection{Economics, incentives, and cost-effectiveness}

A large share of the frontier is economic. Cost--benefit and cost-effectiveness
analyses are repeatedly flagged as future work---for surface amendments and
energy savings \parencite{Wamono2016b}, comprehensive short- and long-term
subsurface-drainage benefits \parencite{Guo2025}, paddy-field cost and
consumption \parencite{Jia2023}, edge-of-field cost-efficiency
\parencite{Mendes2020}, and the opportunity cost of drainage investment together
with policies that motivate improvement
\parencite{Wichelns2002,THOMAS1995,Kladivko2005,Saadat2017,Manik2019,Helmers2012}.
Identifying long-term savings---such as greater fertiliser-nitrogen-use
efficiency---to offset higher installation costs, alongside incentive and
cost-share design, is recommended \parencite{Castellano2019,Sunohara2015}, and
cost reduction is itself a research target through local manufacture and new
low-cost techniques and materials \parencite{Datta2002,Yannopoulos2020}.
Complementary needs include advanced case-by-case cost analyses to promote 4R
implementation \parencite{Liu2021}, more data on the cost-effectiveness and
ecological benefits of conservation practices \parencite{Yuan2022}, and explicit
incorporation of practice costs into future modelling \parencite{McLellan2015}.
Incentive mechanisms warrant study: continued positive incentives and greater
weight on non-excludable benefits in multifunctional landscapes
\parencite{Christianson2014}, cost-share and risk-insurance instruments
\parencite{Jaynes2012,Karlen2005}, and nutrient-trading programmes for
shallow-drain adoption \parencite{Ghane2021b}. Feasibility analyses caution that
drainage-water recycling requires markedly higher, less volatile irrigated
yields and lower adoption costs \parencite{Sellars2022}, motivating work on the
economics of cover crops \parencite{Waring2020}, long-term cumulative yield and
profit of drainage installation \parencite{Mourtzinis2021}, and integrated
biophysical--economic modelling at field scale \parencite{Ferin2023}.

\subsection{Policy, institutions, adoption, and the human dimension}

Biophysical understanding alone is insufficient; governance, capacity, and
farmer engagement are decisive. A repeated recommendation is to shift from
standardised toward flexible, location-specific solutions that reflect
stakeholder preferences, balancing top-down against bottom-up participatory
management and linking farmers' tacit knowledge with researchers' explicit
knowledge \parencite{Ritzema2011,Ritzema2016,Ritzema2023}. Capacity development,
education, and policy change to lower adoption barriers are emphasised
\parencite{Castellano2019,SHIRMOHAMMADI1995,Kesicka2023b}, along with governance
frameworks, inter-agency coordination, dedicated drainage institutions, and
time-bound action plans
\parencite{Singh2019c,Qureshi2008,Gopalakrishnan2007,Ritzema2008b,Valipour2014,Quinn2009},
and awareness, subsidy, micro-finance, and participatory
operation-and-maintenance programmes
\parencite{Satyanarayana2007,Manjunatha2004,BenAissa2013,MCFARLANE1992}.

Adoption and social dimensions are stressed across contexts. Detailed
cost--benefit analyses, joint farmer--researcher experimentation, and capacity
building are prioritised, with adoption ultimately depending on farmers'
long-term cost--benefit perspective \parencite{Okwany2016}; creating know-how and
farmer confidence before project execution is essential \parencite{Ghumman2011};
and key practice-use issues---where to apply BMPs, how to bundle them, and
transferability to commercial farms---remain open \parencite{Ward2018}.
Stakeholder willingness to adopt must inform proposed measures
\parencite{Cuevas2019}, research integrating ecological, pollution-related, and
economic impacts with models for optimal policy design is urged
\parencite{Saikkonen2014}, and incentivising on-farm research while scaling
individual efforts to meso- and macro-scale policy is called for
\parencite{David2015}. Policy reforms sought include legislative reassessment to
remove barriers to natural-water-retention measures \parencite{Kvítek2023},
national support to retain open ditch-and-field networks \parencite{Marja2013},
producer-participation and pilot-programme information for policy optimisation
\parencite{Malone2023}, and integrated packages---reduced plantings, farmer
compensation, and gradual subsidy liberalisation---in salinity contexts
\parencite{Kotb2000}.

\subsection{Integrative landscape-scale frameworks and decision support}

The most explicit long-range vision is numerically explicit landscape-scale
models that couple the processes governing ditch services, optimise the spatial
arrangement of practices according to connectivity, and define multi-service
maximisation strategies informed by stakeholder-valued rankings---prerequisite
to which is better mapping of ditch hydraulic, vegetation, and sorption
properties \parencite{Dollinger2015,Herzon2008}. Landscape-scale approaches that
combine on-farm change with off-site wetland and riparian restoration are
broadly endorsed \parencite{Blann2009,Goeller2019}, and truly integrated
watershed management requires better information on controlling processes such
as the groundwater-head distribution linking topography and geology to the drain
fraction \parencite{Mahmood2023}. Coupled field--watershed models enable
water-table-management scenario analyses not possible with component models used
separately \parencite{Dayyani2010}, frameworks should remain extensible to
additional processes and dynamically calibrated indices
\parencite{Basu2010,Tao2025,Verma2014}, and modelling should increasingly guide
practice placement and stacking
\parencite{Kröger2013,Tomer2015,Lemke2011,McDowell2021,King2022,Roley2016}.

Emerging sensing and data technologies are widely seen as enabling this next
phase. Broad agendas span optimal conditions for detecting drainage patterns and
extracting drain lines, finer-resolution soil moisture, machine-learning mapping
of soil properties, and within-field nitrogen-stress mapping
\parencite{Khanal2020}, alongside spatial frameworks that combine climate and
soil type with geo-referenced farmer-reported data to estimate field-specific
benefits \parencite{Cassman2025}. Decision-support needs include validating
aggregate edge-of-field losses and delivery ratios \parencite{Reid2018},
testing soil-wetness-index tools across soils \parencite{Julien2023}, integrating
water and nutrient management with precision agriculture \parencite{Jung2024},
and coupling downscaled climate information with an understanding of how farmers
incorporate excess-precipitation experience into decisions \parencite{Morton2015}.
Better indices and shared tools are sought---phosphorus and nutrient indices that
represent subsurface drainage, wetland and riparian-buffer removal, and
stream-network connectivity \parencite{Buczko2007,Pärn2012}, extensible
nutrient-credit-trading tools for drainage water management \parencite{Negm2016},
a Soil Vulnerability Index that distinguishes surface from subsurface drainage
\parencite{Baffaut2020}, and coupling with ecosystem-service tools such as INVEST
to assess tradeoffs \parencite{McLellan2015}. Underlying all of these is the data
imperative: better, longer, and more spatially representative datasets---spanning
water and solute fluxes, greenhouse-gas emissions, root and biomass development,
and management history---are the foundational requirement for validating models,
resolving tradeoffs, and delivering site-specific, climate-resilient guidance
\parencite{Ma2025,SKAGGS1994,Deng2011}.
